
\subsection{Point of Diminishing Returns}

The data shows clear diminishing returns beyond Stage 1:

\begin{itemize}[leftmargin=1.5em, itemsep=4pt]
  \item \textbf{Stage 0 → Stage 1:} 
    - Throughput change: -1.3\% (negligible)
    - Quality improvement: +54.1\% (dramatic)
    - Verdict:  Always use Stage 1 over Stage 0 
  
  \item \textbf{Stage 1 → Stage 2:} 
    - Throughput change: -92.8\% (catastrophic)
    - Quality change: -3.5\% (actually worse!)
    - Memory change: -1.5 GB (minimal)
    - Verdict:  Never use Stage 2 
  
  \item \textbf{Stage 2 → Stage 3:} 
    - Throughput change: +28.3\% (improvement over Stage 2)
    - Quality improvement: +18.1\% (better than Stage 2)
    - Memory change: -0.5 GB (minimal)
    - Verdict:  If forced to choose between Stage 2/3, pick Stage 3 
  
  \item \textbf{Stage 1 → Stage 3:} 
    - Throughput change: -90.8\% (unacceptable for most)
    - Quality improvement: +15.2\% (modest)
    - Memory change: -1.0 GB (minimal at this scale)
    - Verdict:  Only for memory-constrained scenarios 
\end{itemize}

 The optimal operating point is clearly Stage 1  — it captures 98.7\% of maximum throughput while achieving 79\% of maximum quality (relative to Stage 3). The extra 15\% quality from Stage 3 costs 10.8× longer training time.


\section{Best Configuration and Conclusions}

\subsection{Best Configuration for This Setup}

Based on the complete empirical evidence,  ZeRO Stage 1 is the optimal configuration  for this 51M parameter model. The supporting data:

\begin{itemize}[leftmargin=1.5em, itemsep=4pt]
  \item \textbf{Throughput:} 2,467,429 tokens/sec — within 1.3\% of Stage 0, and 10-13× higher than Stages 2-3
  
  \item \textbf{Communication overhead:} 0.003\% — identical to Stage 0, dramatically lower than Stages 2 (92.9\%) and 3 (58.9\%)
  
  \item \textbf{Compute efficiency:} 88.2\% — essentially all step time is productive computation
  
  \item \textbf{Model quality:} Final perplexity 2,358 — 54\% better than Stage 0, within 15\% of best (Stage 3 at 1,999)
  
  \item \textbf{Total training time:} 1,066 seconds — only 40 seconds (4\%) longer than Stage 0
  
  \item \textbf{Time to 90\% convergence:} 300 steps — fastest among all stages for reaching good quality
\end{itemize}


\subsection{Conclusions}
\begin{enumerate}
    \item 
\end{enumerate}
\textbf{1. ZeRO Stage 0 must not be used for serious distributed training.}
Without gradient synchronization, each GPU trains independently on its data shard, equivalent to effective batch size of 128 instead of 512. The resulting model achieves 54\% worse perplexity than Stage 1 over the same number of steps (5,133 vs 2,358).

\textbf{2. ZeRO Stage 1 is the optimal stage for this 51M parameter model.}
It delivers near-peak throughput (2.47M tokens/sec, 99\% of Stage 0), 88\% compute efficiency, 0.003\% communication overhead, and strong model quality (2,358 perplexity). The cost of optimizer state partitioning over Stage 0 is a negligible 1.3\% throughput reduction but 54\% quality improvement.

\textbf{3. ZeRO Stages 2 and 3 introduce severe communication bottlenecks at this model scale.}
AllReduce operations consume 92.9\% (Stage 2) and 58.9\% (Stage 3) of total step time, reducing throughput by 90-93\% relative to Stage 1. The root cause is that gradient partitioning forces AllReduce at every microstep rather than only at gradient accumulation boundaries — an 8× increase in communication frequency.

\textbf{4. Expected memory savings from higher ZeRO stages are not realized at this scale. 
Peak memory usage stays within 10.1-11.6 GB across all stages — a span of only 1.5 GB. At 51M parameters, activations and framework overhead dominate memory usage (\approx9.2 GB), and ZeRO only partitions the model state (\approx0.8 GB). The allocated memory (actual usage) drops significantly (1.68 GB → 1.19 GB from Stage 0 to Stage 2), but PyTorch's caching allocator masks this benefit in peak memory measurements.

\textbf{5. Communication patterns differ qualitatively between stages, not just quantitatively.}
\begin{itemize}
    \item Stages 0-1: AllReduce is infrequent (every 8 microsteps) and bursty — spikes of 1,300ms followed by 7 microsteps of 0ms
    \item Stages 2-3: AllReduce occurs at every single microstep as constant, persistent overhead — flat line at 10,600ms or 5,300ms
\end{itemize}

This qualitative difference explains why the jump from Stage 1 to Stage 2 results in a 92.9\% communication overhead — the communication pattern completely changes, not just scales up.

\textbf{6. Stage 3's forward pass cost is uniquely elevated due to parameter reconstruction.}
Parameter AllGather operations in Stage 3 add approximately 2,870ms to each forward pass relative to other stages (3,187ms vs 310ms). This cost is entirely absent in Stages 0, 1, and 2. For small models, this overhead dominates; for large models (>10B parameters), the AllGather cost can be amortized over larger parameter sizes.

\textbf{7. ZeRO Stages 2-3 become appropriate only when the model cannot fit in GPU memory at lower stages.}
In scenarios with much larger models (>1B parameters), per-GPU memory pressure would make Stages 0-1 infeasible, and Stages 2-3 would be necessary. In such scenarios, communication overhead would need to be managed through:
- Gradient checkpointing (reduces activation memory by 50-70\%)
- Activation partitioning (ZeRO-3's activation offloading)
- Larger global batch sizes (amortize communication across more computation)
- High-bandwidth interconnects (NVLink: 600 GB/s vs PCIe 4.0: 32 GB/s)
- Communication-computation overlap (overlap AllReduce with backward compute)







 Why not Stage 0?  Despite 1.3\% higher throughput, Stage 0's lack of gradient synchronization produces 54\% worse perplexity (5,133 vs 2,358). The throughput advantage is meaningless when the model learns so poorly.

 Why not Stage 2?  Worse in every metric: 92.8\% lower throughput, 3.5\% worse quality, and the highest communication overhead. Stage 2 has no redeeming qualities.

 Why not Stage 3?  While achieving best final quality (1,999 perplexity), the 90.8\% throughput reduction (2.47M → 0.23M tok/s) means Stage 3 takes 4.3× longer to train for only 15\% quality improvement. For most applications, the extra 15\% quality isn't worth 4.3× longer training.


\subsection{Final Recommendation}

For the vast majority of distributed training scenarios with models that fit in GPU memory (up to approximately 1B parameters on 24GB GPUs),  ZeRO Stage 1 is the optimal choice . It provides:
\begin{itemize}
    \item 98\% of maximum possible throughput
    \item 54\% better model quality than no synchronization
    \item 4 \times faster training than Stage 3 for equivalent quality
    \item Negligible communication overhead  (0.003\% vs 92\% in Stage 2)
\end{itemize}

Only when model size exceeds GPU memory capacity should Stage 3 be considered, and Stage 2 should never be used. The empirical data is unambiguous: Stage 2 is dominated by Stage 3 in every meaningful metric.

\subsection{Future Work}

This analysis suggests several directions for future investigation:
\begin{enumerate}
    \item Scale to larger models (1B, 10B parameters) to observe where ZeRO stages 2-3 become memory-advantageous
    \item Test with higher-bandwidth interconnects (NVLink, InfiniBand) to reduce communication overhead
    \item Implement communication-computation overlap to hide AllReduce latency
    \item Explore hybrid strategies (ZeRO-1 for some layers, ZeRO-3 for others)
    \item Measure impact of gradient accumulation steps
    \item Analyze activation memory separately from model state memory
    \item Profile with larger batch sizes to increase compute/communication ratio
\end{enumerate}


\appendix
\section{Chart-by-Chart Observation Reference}

This reference table provides quick look-up of key observations for each figure:

\renewcommand{\arraystretch}{1.3}
\begin{longtable}{L{3.5cm} L{10.5cm}}
\toprule
\rowcolor{primaryblue}
\textcolor{white}{\textbf{Chart}} & \textcolor{white}{\textbf{Key Observation}} \\
\midrule
\endfirsthead
\toprule
\rowcolor{primaryblue}
\textcolor{white}{\textbf{Chart}} & \textcolor{white}{\textbf{Key Observation}} \\
\midrule
\endhead
\bottomrule
\endfoot
\rowcolor{lightgray}
Figure 1: Loss Convergence & Stage 1 drops fastest initially; Stage 3 achieves lowest final loss; Stage 0 worst throughout. \\
Figure 2: Loss Reduction & Stage 3 highest total reduction (2.93); Stage 1 reaches 90\% reduction fastest (300 steps). \\
\rowcolor{lightgray}
Figure 3: Throughput Bar & Stages 0-1 at ~2.5M tok/s; Stages 2-3 at ~0.2M tok/s (10-12× drop). \\
Figure 4: Iteration Time & Stage 0 fastest (2.14s/iter); Stage 2 slowest (11.56s); Stage 3 intermediate (9.14s). \\
\rowcolor{lightgray}
Figure 5: Memory Bar & All stages 10.1-11.6 GB; no significant memory gradient; Stage 1 highest (overhead). \\
Figure 6: Memory-Throughput Tradeoff & Stages 0-1 dominate top; Stages 2-3 lower-left; no stage optimal on both axes. \\
\rowcolor{lightgray}
Figure 7: GPU Utilization & All stages 97.7-99.9\%; high utilization but Stage 2 wastes 93\% on communication. \\
Figure 8: Comm Overhead Bar & Stages 0-1 0.003\%; Stage 2 92.9\%; Stage 3 58.9\% — the performance cliff. \\
\rowcolor{lightgray}
Figure 9: Comm Overhead/Microstep & Stage 0 →0\%; Stage 1 →6\%; Stage 2 flat at →48\%; Stage 3 flat at →37\%. \\
Figure 10: Forward Pass/Microstep & Stages 0-2 ~300ms; Stage 3 ~3,200ms (10.7× slower due to AllGather). \\
\rowcolor{lightgray}
Figure 11: Backward Pass/Microstep & Stages 0-1 ~1,850ms; Stage 2 ~11,000ms; Stage 3 ~5,900ms. \\
Figure 12: AllReduce/Microstep & Stage 0 near-zero; Stage 1 oscillating; Stage 2 flat ~10,600ms; Stage 3 ~5,300ms. \\
\rowcolor{lightgray}
Figure 13: Total Microstep Time & Stages 0-1 ~3.5s/8 microsteps; Stage 3 ~14.5s; Stage 2 ~22s. \\
Figure 14: Compute Efficiency & Stages 0-1 88-94\%; Stage 3 63\%; Stage 2 worst at 52\% (communication dominates). \\
\rowcolor{lightgray}
Figure 15: Stacked Components & Stages 0-1 sparse AllReduce spikes; Stages 2-3 constant dense communication. \\
Figure 16: Component Box Plots & Tight distributions for all stages — stable, reproducible behavior. \\
\rowcolor{lightgray}
Figure 17: Avg Components with Error Bars & Large error bars for Stages 0-1 (bursty); tight bars for Stages 2-3 (constant). \\
Figure 18: Correlation Heatmap & bwd↔allreduce=1.00; fwd↔stage=0.77; bwd↔stage=0.52. \\
\rowcolor{lightgray}
Figure 19: Radar Chart & Stages 0-1 fill pentagon; Stages 2-3 collapsed polygon — not competitive. \\
Figure 20: Final Perplexity & Stage 0 worst (5,133); Stage 3 best (1,999); Stage 1 good (2,358). \\
\rowcolor{lightgray}
Figure 21: Perplexity Improvement & Stages 1/2/3 improve 54.1\%, 52.5\%, 61.1\% over Stage 0 — dramatic quality gains. \\
Figure 22: Perplexity Over Steps & Stage 0 slower convergence; Stages 1-3 similar; Stage 3 slightly best final. \\
\rowcolor{lightgray}
Figure 23: Perplexity Subplots & Stage 0 starts higher; Stage 2 shows instability bump steps 100-130; Stage 3 best end. \\
Figure 24: Memory-Perplexity Tradeoff & No stage optimal on both; Stage 3 best quality; Stage 2 best memory. \\
\end{longtable}

\end{document}
